\section{Problem Formulation}
\label{sec:problem}

% 证据：prior/.../多型号多地/mtms.py（步长）、tools_model.py（工具占用）、freshlb.py（开班时刻）；研究笔记/01 §2、07。
% 用户决定：多站点单基地（10-27）、提前大修只在启动窗口内（10-29）、单目标（10-30）。
% 写法（11-04 用户意见）：实例数值放实验一节；问题直接在格上定义，小时换算成格放实验一节。
% 11-04 第二次：仿照 RDD 稿第三节和 Bertsimas–Sethuraman 2002 的写法，问题用文字讲清，只保留目标式；
%   时间原点取"所有机型都能出发的第一格"（原 c_0），于是站点 s 从 λ_s 开始值守。

\begin{table}[t]
  \centering
  \caption{Notation}
  \label{tab:notation}
  \footnotesize
  \begin{tabular}{@{}ll@{}}
    \toprule
    Symbol & Meaning \\
    \midrule
    $\mathcal{S}$, $K_s$ & sites; robots needed on duty at site $s$ \\
    $\lambda_s$ & travel time from the depot to site $s$ \\
    $\mathcal{K}$, $\mathcal{E}_s$ & robot types; types allowed at site $s$ \\
    $d_k$ & duty length of type $k$ \\
    $L_k$ & duties in one maintenance cycle of type $k$ \\
    $\mu_{k,i,s}$ & time a phase-$i$ duty at site $s$ occupies a type-$k$ robot \\
    $\mathcal{J}$, $\mathcal{C}_k$, $q_k$ & tool types; types compatible with $k$; tools per duty \\
    $\alpha_{k,s}$, $\beta_{k,s,j}$ & tool busy time before and after the duty start \\
    $W$ & length of the start-up window \\
    $N_k$, $M_j$ & numbers of type-$k$ robots and type-$j$ tools \\
    \bottomrule
  \end{tabular}
\end{table}

A new fleet of robots must keep several sites staffed without end.
We seek the smallest fleet of robots and tools that can do so.
Time is divided into slots of equal length, and all durations are integer numbers of slots.
Table~\ref{tab:notation} lists the notation.

\textbf{Sites.}
The robots start from one depot.
Site $s$ is $\lambda_s$ slots away and needs exactly $K_s$ robots on duty in every slot; only robots of types in $\mathcal{E}_s$ may serve there.

\textbf{Robots.}
A robot of type $k$ serves duties of $d_k$ slots and is maintained after every duty.
The maintenance follows a cycle of $L_k$ duties: a short maintenance after each of the first $L_k - 1$ duties, and an overhaul after the last.
The position of a duty in this cycle is its \emph{phase} $i \in \{0, \dots, L_k - 1\}$.
A phase-$i$ duty at site $s$ occupies the robot for $\mu_{k,i,s}$ slots from its start: the duty, the trip back, the maintenance and the preparation of the next duty.
The robot then travels to its next site, possibly after waiting, and serves its next duty in phase $i + 1 \bmod L_k$.

\textbf{Tools.}
A duty of type $k$ needs $q_k$ tools whose types are in $\mathcal{C}_k$.
A type-$j$ tool is busy from $\alpha_{k,s}$ slots before the duty starts, for installation and the trip, until $\beta_{k,s,j}$ slots after it starts, when it has been removed and serviced.
A tool serves one duty at a time and needs no maintenance cycle.

\textbf{Start-up.}
At slot $0$ every robot is new and in phase $0$, which is the state right after an overhaul.
Site $s$ must be staffed from slot $\lambda_s$, the first slot a robot can reach it.
In the start-up window $[0, W)$ a robot may take the overhaul after any duty and return to phase $0$; after the window it follows its cycle.

A \emph{schedule} gives every robot an infinite sequence of duties, waits and early overhauls, and gives every duty its tools.
It is \emph{feasible} if no robot or tool is busy with two duties at once and every site is staffed as required.
Multi-domain rotation rostering (MDRR) asks for a feasible schedule that minimizes
\begin{equation}
  \textstyle\sum_{k \in \mathcal{K}} N_k + \sum_{j \in \mathcal{J}} M_j ,
  \label{eq:objective}
\end{equation}
the number of robots and tools; a domain is a site.

The horizon is infinite because a fleet is bought once and serves for years.
A schedule that is feasible over a finite horizon can run short later; Section~\ref{sec:experiments} shows how often this happens when a search optimizes a finite horizon.
